% ─────────────────────────────────────────────────────────────────────────────
\section{Introduction and Scope}
% ─────────────────────────────────────────────────────────────────────────────

This document records the mathematical foundations and algorithmic design decisions
behind \catena, a Python library for simple continued fractions.
It is intended as a living document: new sections will be added as the library
grows to cover further aspects of the theory.

The current implementation (\texttt{v0.4.0}) covers finite SCFs over the rationals,
infinite generative SCFs, and periodic SCFs for quadratic irrationals.

\subsection{What continued fractions offer}

A real number can be represented in several ways: as a floating-point approximation,
as an exact rational via \texttt{fractions.Fraction}, as a fixed-precision decimal,
or as a continued fraction.  Each representation makes different trade-offs.

Floats are fast and ubiquitous, but they introduce rounding error at every operation.
A float is not a rational number; it is the nearest representable binary fraction,
and the gap between the two is invisible unless you care to look for it.

\texttt{fractions.Fraction} is exact and supports arbitrary precision, but it
accumulates large numerators and denominators quickly.  Adding two fractions
$p_1/q_1 + p_2/q_2$ yields $(p_1 q_2 + p_2 q_1)/(q_1 q_2)$; in the worst case
(coprime denominators) the denominator grows from $n$ digits to $2n$ digits before
any reduction, and every operation incurs a GCD computation whose cost grows with
digit count.

Continued fractions offer a different kind of exactness.  The partial quotients
$[a_0;\, a_1, a_2, \ldots]$ are small integers — typically single digits for
naturally occurring constants — so the representation is compact.  The $n$-th
convergent $p_n/q_n$ is automatically in lowest terms by the determinant identity
$p_n q_{n-1} - p_{n-1} q_n = \pm 1$, with no GCD step required.  Convergents
also have an optimality property: $p_n/q_n$ is the best rational approximation
to the true value among all fractions with denominator $\leq q_n$.

For irrational numbers, the continued fraction is the natural representation.
$\sqrt{2} = [1;\overline{2}]$, $\varphi = [1;\overline{1}]$, and more generally
every quadratic irrational has an eventually periodic expansion.  Floats and
decimals can only approximate; the SCF records the structure exactly.

These properties make SCFs useful for problems where rational approximation
quality matters — gear-ratio design, modular arithmetic, Pell's equation, and
the analysis of approximation algorithms, among others.  \catena is built around
this perspective: the library works with the SCF itself as the primary object,
deriving convergents and numeric approximations from it on demand.

\subsection{Scope}

\catena provides three main abstractions:

\begin{itemize}[nosep]
  \item \texttt{SimpleContinuedFraction} — an infinite SCF $[a_0;\,a_1,a_2,\ldots]$
        driven by an arbitrary callable generator.  Convergents are memoised as they
        are requested; the cache is shared when the same tail is viewed with a different
        integer part (e.g.\ after an integer shift).
  \item \texttt{FiniteSimpleContinuedFraction} — a finite SCF
        $[a_0;\,a_1,\ldots,a_k]$ over $\Q$, constructible from a rational, float,
        or decimal string.  Supports exact rational arithmetic and multiplicative inverse.
  \item \texttt{PeriodicSimpleContinuedFraction} — an eventually-periodic SCF
        $[a_0;\,a_1,\ldots,a_m,\,\overline{b_1,\ldots,b_r}]$ representing a quadratic
        irrational $(P + \sqrt{D})/Q$.  Provides algebraic operations: conjugate,
        inverse, negation, and exact decimal evaluation.
\end{itemize}

\begin{remark}
  \texttt{PeriodicSimpleContinuedFraction} is listed here for completeness.
  Its full mathematical treatment — Lagrange's theorem on quadratic irrationals,
  the Pell equation, and the algebraic operations
  (\texttt{conjugate}, \texttt{inverse}, \texttt{quadratic\_coefficients}) —
  is deferred to a forthcoming section.
\end{remark}

The mathematical content of this document is therefore structured around these
three cases.  It is not intended to be an exhaustive treatment of continued fraction
theory; for that, see \cite{khinchin1964} and \cite{rockett1992}.  The goal is to
give enough background to justify the algorithms and design choices in the library.

\subsection{Conventions}
      
Throughout this document:
\begin{itemize}[nosep]
      \item $\N$ denotes the natural numbers (including zero).
      \item For any $n \in \N$, we implicitly assume the set-theoretic construction $n = \{0, 1, \ldots, n-1\}$.
      \item $\Z$, $\Q$, $\R$ denote the integers, rationals, and reals respectively.
      \item Likewise, $\Zp$, $\Qp$, $\Rp$ denote the positive (non-zero) integers, rationals, and reals respectively.
      \item $[a_0;\, a_1, a_2, \ldots, a_n]$ denotes a finite SCF with $a_0 \in \Z$ and
            $a_k \in \Zp$ for $k \geq 1$.
      \item $c_k = (p_k, q_k)$ denotes the $k$-th convergent numerator--denominator pair, with
            $q_k \in \Zp$ by convention.
      \item The function $\{\cdot\}: \R \to \R$ denotes the fractional part, so $\{x\} = x - \lfloor x \rfloor$.
\end{itemize}

\begin{remark}
  We follow the set-theoretic convention $\N = \{0, 1, 2, \ldots\}$ including zero.
  The positive integers $\Zp = \{1, 2, 3, \ldots\}$ exclude zero.
  In particular, all body partial quotients $a_k$ (for $k \geq 1$) lie in $\Zp$,
  while the integer part $a_0 \in \Z$ may be zero or negative.
\end{remark}